\section{Introducción}

El despliegue de una red de telefonía móvil exige decidir en cuáles de un conjunto de ubicaciones candidatas se instalarán antenas para que todos los usuarios queden dentro del área de cobertura de, al menos, una de ellas. Cada ubicación tiene un costo de instalación distinto y cubre a un subconjunto particular de clientes. El objetivo de la empresa es doble: no dejar a ningún cliente sin servicio y hacerlo con el menor costo total posible.

Este problema es una instancia clásica del \textit{problema de cobertura de conjuntos} (\textit{Set Cover Problem}, SCP), uno de los 21 problemas NP-completos de Karp \cite{karp72}. Aparece, además de en telecomunicaciones, en la asignación de personal, la localización de servicios de emergencia, la selección de pruebas de software y el diseño de circuitos. Aunque su enunciado es sencillo, su dificultad computacional crece con el tamaño de la instancia: no se conoce un algoritmo que lo resuelva de forma eficiente en el caso general \cite{garey79}, y resulta difícil incluso de aproximar mejor que un factor logarítmico \cite{feige98}.

En este trabajo se considera una instancia con 500 clientes y 500 ubicaciones candidatas, cuyos subconjuntos de cobertura se entregan en una matriz binaria y cuyos costos se entregan en un archivo aparte. El trabajo plantea el modelo de programación lineal entera binaria del Set Cover y lo resuelve mediante un Algoritmo Genético (AG) diseñado para este problema. Además, para disponer de un punto de referencia, se resuelve el mismo modelo con un solver exacto y se comparan ambos enfoques en calidad de la solución y tiempo de cómputo.

\subsection{Objetivos}
\textbf{Objetivo general.} Determinar el conjunto de antenas de menor costo que cubra a los 500 clientes, modelando el problema como Set Cover y resolviéndolo con un Algoritmo Genético.

\textbf{Objetivos específicos.}
\begin{itemize}
\item Formular el modelo de programación lineal entera binaria del problema, definiendo variables, función objetivo y restricciones.
\item Diseñar e implementar un Algoritmo Genético con operadores de selección, cruce, mutación y reparación adecuados para el Set Cover.
\item Resolver el modelo con un solver exacto y usar su resultado como referencia.
\item Comparar ambos métodos en calidad de solución, tiempo de cómputo y comportamiento, y analizar el efecto de los parámetros del AG.
\end{itemize}

\subsection{Estructura del documento}
La sección 2 presenta el marco teórico. La sección 3 describe la metodología: datos, modelo matemático, método exacto y Algoritmo Genético. La sección 4 reporta y analiza los resultados. La sección 5 presenta las conclusiones.
